% !TEX root = ../main.tex
%% Corrigés — Chapitre 12 — Espaces vectoriels et applications linéaires
%% Résultats vérifiés par verifications/chapitre_12.py.

\corrigesChapitre{12}{Espaces vectoriels et applications linéaires}
\corrigesSouples

\begin{blocLivre}{cor-12-1}
\begin{corrige}{1}{Sous-espace ou pas ?}
\demarche{Trois réflexes : le vecteur nul est-il dedans ? la somme reste-t-elle
dedans ? un multiple reste-t-il dedans ?}
\begin{questions}
  \item $\ens{x+2y=0}$ : \res{\text{oui}} (contient $(0,0)$, stable par
    combinaison : c'est le noyau de la forme linéaire $(x,y)\mapsto x+2y$).
    $\ens{x+y=1}$ : \res{\text{non}}, $(0,0)$ n'y est pas.
  \item \res{\text{Non}} : $(1,0)$ et $(0,1)$ y sont, mais leur somme
    $(1,1)$ non — la réunion des deux axes n'est pas stable.
  \item $\ens{f(1)=0}$ : \res{\text{oui}} (noyau de l'évaluation en $1$).
    $\ens{f(1)=1}$ : \res{\text{non}} (la fonction nulle n'y est pas).
  \item \res{\text{Non}} : le polynôme nul manque ($\deg0=-\infty$), et
    $X^{3}$ et $-X^{3}+X$ sont de degré $3$ alors que leur somme $X$ est de
    degré $1$.
  \item \res{\text{Oui}} : $0$ est bornée, et si $|u_n|\leqslant M$,
    $|v_n|\leqslant M'$, alors $|u_n+\lambda v_n|\leqslant M+|\lambda|M'$.
\end{questions}
\end{corrige}
\end{blocLivre}

\begin{blocLivre}{cor-12-2}
\begin{corrige}{2}{Dans un Vect, ou pas ?}
\demarche{Appartenir à un Vect, c'est être solution d'un système.}
\begin{questions}
  \item $(1,2,3)=au+bv$ donne $a=1$, $a+b=2$, $a+2b=3$ : solution
    $(a,b)=(1,1)$, donc $(1,2,3)=u+v$ : \res{\text{oui}}. Pour $(1,0,0)$ :
    $a=1$, $a+b=0$ donc $b=-1$, mais $a+2b=-1\neq0$ : \res{\text{non}}.
  \item $(x,y,z)=au+bv$ impose $a=x$, $b=y-x$, puis
    $z=a+2b=2y-x$ : le plan a pour équation
    $\res{x-2y+z=0}$ (vérification : $u$ et $v$ la satisfont).
  \item $(1,2,c)$ : $1-4+c=0$, soit $\res{c=3}$ — cohérent avec a).
\end{questions}
\end{corrige}
\end{blocLivre}

\begin{blocLivre}{cor-12-3}
\begin{corrige}{3}{Libres ou liées ?}
\demarche{Une combinaison nulle, un système — et pour les fonctions, des
évaluations bien choisies.}
\begin{questions}
  \item $a(1,0,1)+b(1,1,0)+c(0,1,1)=0$ donne $a+b=0$, $b+c=0$, $a+c=0$,
    d'où $a=-b=c=-a$ : $a=b=c=0$, famille \res{\text{libre}}.
  \item $(2,4,6)=2\,(1,2,3)$ : famille \res{\text{liée}} (colinéarité).
  \item $(1,1,0)-(0,1,1)-(1,0,-1)=(0,0,0)$ : famille \res{\text{liée}},
    avec la relation $v_1-v_2-v_3=0$.
  \item $(1,X,\dots,X^{n})$ : degrés échelonnés, donc libre (un polynôme
    nul a tous ses coefficients nuls). Pour $a\cos+b\sin=0$ : en $0$,
    $a=0$ ; en $\frac{\pi}{2}$, $b=0$ : famille \res{\text{libre}}.
\end{questions}
\end{corrige}
\end{blocLivre}

\begin{blocLivre}{cor-12-4}
\begin{corrige}{4}{Une base échelonnée de $\R_3[X]$}
\demarche{Degrés échelonnés = liberté gratuite ; le bon compte de vecteurs
fait le reste.}
\begin{questions}
  \item Degrés $0,1,2,3$ deux à deux distincts : famille libre de $4$
    vecteurs dans $\R_3[X]$, qui est de dimension $4$ : c'est une
    \res{\text{base}}.
  \item $X^{3}=\bigl((X-1)+1\bigr)^{3}
    =(X-1)^{3}+3(X-1)^{2}+3(X-1)+1$ (binôme) : coordonnées
    $\res{(1,3,3,1)}$ dans $\mathcal{B}$.
  \item La formule de Taylor $P=\sum_{k=0}^{3}\dfrac{P^{(k)}(1)}{k!}(X-1)^{k}$
    (chapitre 10) : les coordonnées dans $\mathcal{B}$ sont les
    $\dfrac{P^{(k)}(1)}{k!}$.
\end{questions}
\end{corrige}
\end{blocLivre}

\begin{blocLivre}{cor-12-5}
\begin{corrige}{5}{Noyau et image d'une application linéaire}
\demarche{Le noyau se calcule par un système, l'image par les colonnes —
et le théorème du rang recolle tout.}
\begin{questions}
  \item Chaque composante est une combinaison linéaire des coordonnées :
    $u(X+\lambda Y)=u(X)+\lambda u(Y)$ s'ensuit coordonnée par coordonnée.
  \item $x+y-z=0$ et $2x-y=0$ donnent $y=2x$ puis $z=3x$ :
    $\operatorname{Ker}u=\res{\operatorname{Vect}\bigl((1,2,3)\bigr)}$,
    de dimension $1$.
  \item $\operatorname{Im}u=\operatorname{Vect}\bigl(u(e_1),u(e_2),u(e_3)\bigr)
    =\operatorname{Vect}\bigl((1,2),(1,-1),(-1,0)\bigr)$ : les deux premiers
    vecteurs ne sont pas colinéaires, donc
    $\operatorname{Im}u=\R^{2}$ et $\operatorname{rg}u=\res{2}$. Et
    $\dim\R^{3}=3=1+2$ : le théorème du rang est confirmé.
\end{questions}
\end{corrige}
\end{blocLivre}

\begin{blocLivre}{cor-12-6}
\begin{corrige}{6}{Linéaire ou pas ?}
\demarche{Un seul contre-exemple suffit pour le non ; pour le oui, la
définition en une ligne.}
\begin{questions}
  \item $f$ : \res{\text{oui}} (composantes linéaires en $(x,y)$).
    $g$ : \res{\text{non}}, $g(0,0)=(1,0)\neq(0,0)$ — une application
    linéaire envoie toujours $0$ sur $0$.
  \item $h$ : \res{\text{non}} : $h(2,2)=4\neq2=2\,h(1,1)$.
  \item $(P+\lambda Q)'=P'+\lambda Q'$ : \res{\text{oui}} ; de même
    $(P+\lambda Q)(X+1)=P(X+1)+\lambda Q(X+1)$ : \res{\text{oui}}.
  \item $(f+\lambda g)(2)=f(2)+\lambda g(2)$ : \res{\text{oui}}, c'est une
    forme linéaire sur $\R^{\R}$.
\end{questions}
\end{corrige}
\end{blocLivre}

\begin{blocLivre}{cor-12-7}
\begin{corrige}{7}{Un projecteur à la loupe}
\demarche{Un projecteur se lit sur son image (les points fixes) et son noyau
(la direction de projection).}
\begin{questions}
  \item Linéarité claire ; en posant $m=\dfrac{x+y}{2}$,
    $p(p(x,y))=p(m,m)=\Bigl(\dfrac{m+m}{2},\dfrac{m+m}{2}\Bigr)=(m,m)$ :
    $p\circ p=p$.
  \item $\operatorname{Im}p=\operatorname{Vect}\bigl((1,1)\bigr)$ (la
    diagonale $\Delta$, dont les points sont fixes) et
    $\operatorname{Ker}p=\operatorname{Vect}\bigl((1,-1)\bigr)$
    ($x+y=0$) : $p$ est la \res{\text{projection sur }\Delta\text{
    parallèlement à l'antidiagonale}}.
  \item $s(x,y)=2p(x,y)-(x,y)=\res{(y,\,x)}$ : l'échange des coordonnées !
    $s(s(x,y))=s(y,x)=(x,y)$ : $s^{2}=\mathrm{id}$, et $s$ est la symétrie
    par rapport à $\Delta$ parallèlement à l'antidiagonale.
\end{questions}
\medskip
\begin{center}
\begin{tikzpicture}[scale=0.9,line cap=round]
  \draw[otMuted,line width=0.5pt,-{Stealth[length=1.6mm]}] (-1.9,0) -- (2.3,0);
  \draw[otMuted,line width=0.5pt,-{Stealth[length=1.6mm]}] (0,-1.5) -- (0,2.1);
  \draw[otNavy,line width=1pt] (-1.5,-1.5) -- (2.0,2.0);
  \draw[otGreenDark,line width=0.9pt] (1.4,-1.4) -- (-1.4,1.4);
  \draw[otInk,line width=1pt,-{Stealth[length=2mm]}] (0,0) -- (1.8,0.6);
  \draw[otMuted,line width=0.6pt,dash pattern=on 2pt off 1.5pt] (1.8,0.6) -- (1.2,1.2);
  \draw[otNavy,line width=0.9pt,-{Stealth[length=2mm]}] (0,0) -- (1.2,1.2);
  \draw[otMuted,line width=0.6pt,dash pattern=on 2pt off 1.5pt] (1.8,0.6) -- (0.6,1.8);
  \draw[otGreenDark,line width=0.9pt,-{Stealth[length=2mm]}] (0,0) -- (0.6,1.8);
  \node[font=\scriptsize,text=otInk,anchor=west] at (1.85,0.6) {$(x,y)$};
  \node[font=\scriptsize,text=otNavy,anchor=north west] at (1.25,1.1) {$p(x,y)$};
  \node[font=\scriptsize,text=otGreenDark,anchor=south west] at (0.62,1.8) {$s(x,y)=(y,x)$};
  \node[font=\scriptsize,text=otNavy,anchor=south] at (2.0,2.02) {$\Delta$};
\end{tikzpicture}\\[0.5mm]
\legende{Projection sur la diagonale $\Delta$ et symétrie associée :
$s$ échange les coordonnées.}
\end{center}
\end{corrige}
\end{blocLivre}

\begin{blocLivre}{cor-12-8}
\begin{corrige}{8}{Rangs et dimensions}
\demarche{Le rang, c'est le nombre de vecteurs réellement utiles.}
\begin{questions}
  \item D'après l'exercice 3 c), la famille est liée mais
    $\bigl((1,1,0),(0,1,1)\bigr)$ est libre : rang $\res{2}$. Les trois
    vecteurs $(1,2),(2,4),(3,6)$ sont colinéaires à $(1,2)$ : rang
    $\res{1}$.
  \item Hyperplan de $\R^{3}$ (noyau de $(x,y,z)\mapsto x+y+z$) :
    dimension $\res{2}$, base $\bigl((1,-1,0),(1,0,-1)\bigr)$ (deux
    vecteurs non colinéaires de somme des coordonnées nulle).
  \item $a(1+X)+b(1-X)+cX^{2}=0$ donne $c=0$, $a+b=0$ et $a-b=0$, d'où
    $a=b=c=0$ : famille libre de $3$ vecteurs dans $\R_2[X]$ de dimension
    $3$ : une \res{\text{base}}.
\end{questions}
\end{corrige}
\end{blocLivre}

\begin{blocLivre}{cor-12-9}
\begin{corrige}{9}{Des espaces de solutions}
\demarche{Dimension 2 = deux conditions initiales : une base à deux éléments
décrit toutes les solutions.}
\begin{questions}
  \item C'est le noyau de l'application linéaire $y\mapsto y''+y$ : un
    sous-espace vectoriel. Si $a\cos+b\sin=0$, évaluer en $0$ et
    $\frac{\pi}{2}$ donne $a=b=0$ : famille libre. En dimension $2$, elle
    est donc une base : \emph{toute} solution s'écrit
    $\res{a\cos+b\sin}$.
  \item $\Phi:u\mapsto(u_0,u_1)$ est linéaire ; injective, car $u_0=u_1=0$
    propage $u_n=0$ par récurrence ; surjective, car tout couple de
    conditions initiales définit une suite de $E$ par la récurrence. Donc
    $\Phi$ est un isomorphisme et $\dim E=\res{2}$.
  \item $(r^{n})\in E$ pour $r\neq0$ si et seulement si $r^{2}=r+1$, dont
    les racines sont $\varphi$ et $\psi$. Si
    $a\varphi^{n}+b\psi^{n}=0$ pour tout $n$, les rangs $n=0,1$ donnent
    $a+b=0$ et $a\varphi+b\psi=0$, système de déterminant
    $\psi-\varphi=-\sqrt5\neq0$ : $a=b=0$. Famille libre de deux vecteurs
    en dimension $2$ : une base. Pour Fibonacci ($F_0=0$, $F_1=1$) :
    $a+b=0$ et $a\varphi+b\psi=1$ donnent $a=\dfrac{1}{\sqrt5}=-b$, d'où
    \[
    \res{F_n=\frac{\varphi^{n}-\psi^{n}}{\sqrt5}}
    \]
    — la formule de Binet (chapitre 5), retrouvée par pure algèbre
    linéaire.
\end{questions}
\end{corrige}
\end{blocLivre}

\begin{blocLivre}{cor-12-10}
\begin{corrige}{10}{Supplémentaires : deux grands classiques}
\demarche{En dimension finie : intersection nulle + compte des dimensions.
En dimension infinie : analyse-synthèse.}
\begin{questions}
  \item $F\cap G$ : $(t,t,t)$ avec $3t=0$, donc $\ens{0}$ ; et
    $\dim F+\dim G=2+1=3=\dim\R^{3}$ : $\R^{3}=F\oplus G$. Pour $(1,2,3)$ :
    chercher $g=\lambda(1,1,1)$ ; la somme des coordonnées de $f$ étant
    nulle, $6=3\lambda$, donc $\lambda=2$ :
    \[
    (1,2,3)=\res{(-1,0,1)+(2,2,2)}.
    \]
  \item \emph{Analyse} : si $f=p+i$ avec $p$ paire, $i$ impaire, alors
    $f(-x)=p(x)-i(x)$, d'où $p(x)=\dfrac{f(x)+f(-x)}{2}$ et
    $i(x)=\dfrac{f(x)-f(-x)}{2}$ : unicité. \emph{Synthèse} : ces formules
    définissent bien une paire et une impaire de somme $f$. Pour
    l'exponentielle : $p=\operatorname{ch}$ et $i=\operatorname{sh}$ —
    $\res{\mathrm{e}^{x}=\operatorname{ch}x+\operatorname{sh}x}$.
\end{questions}
\end{corrige}
\end{blocLivre}

\begin{blocLivre}{cor-12-11}
\begin{corrige}{11}{La formule de Grassmann en action}
\demarche{Grassmann transforme une question géométrique en arithmétique des
dimensions.}
\begin{questions}
  \item $F+G$ contient strictement $F$ (sinon $G\subseteq F$ et, par
    égalité des dimensions, $G=F$), donc $\dim(F+G)=3$ et
    \[
    \dim(F\cap G)=2+2-3=1 :
    \]
    l'intersection est une \res{\text{droite vectorielle}}.
  \item Un vecteur de $G$ s'écrit $a(0,0,1)+b(1,1,1)=(b,b,a+b)$ ; il est
    dans $F$ (troisième coordonnée nulle) si et seulement si $a+b=0$, et
    vaut alors $b\,(1,1,0)$ :
    $F\cap G=\res{\operatorname{Vect}\bigl((1,1,0)\bigr)}$.
  \item Non : $\dim(F\cap G)=3+3-\dim(F+G)\geqslant6-5=1$ : dans $\R^{5}$,
    deux sous-espaces de dimension $3$ se coupent \emph{toujours} hors de
    l'origine.
\end{questions}
\end{corrige}
\end{blocLivre}

\begin{blocLivre}{cor-12-12}
\begin{corrige}{12}{Base incomplète, base extraite}
\demarche{Compléter : ajouter un vecteur hors du Vect. Extraire : jeter les
vecteurs superflus.}
\begin{questions}
  \item $(1,0,0)\notin\operatorname{Vect}\bigl((1,1,0),(0,1,1)\bigr)$ :
    sinon $(1,0,0)=a(1,1,0)+b(0,1,1)$ avec $a=1$, $a+b=0$, $b=0$ :
    impossible. La famille $\res{\bigl((1,1,0),(0,1,1),(1,0,0)\bigr)}$ est
    libre, donc base de $\R^{3}$.
  \item $(2,0)=2(1,0)$ et $(1,1)=(1,0)+(0,1)$ sont superflus :
    $\res{\bigl((1,0),(0,1)\bigr)}$ est extraite et génératrice, donc base.
  \item La famille $\bigl(1+X+X^{2},\,X,\,1\bigr)$ est libre (ses
    coordonnées dans $(1,X,X^{2})$ forment un système triangulaire
    inversible), donc $\R_2[X]=\operatorname{Vect}(1+X+X^{2},X)\oplus
    \res{\operatorname{Vect}(1)}$ — tout autre choix hors du plan convient
    aussi : un supplémentaire n'est jamais unique.
\end{questions}
\end{corrige}
\end{blocLivre}

\begin{blocLivre}{cor-12-13}
\begin{corrige}{13}{Détermination d'une application linéaire}
\demarche{Prescrire l'image d'une base détermine tout : il n'y a plus qu'à
décomposer.}
\begin{questions}
  \item $\bigl((1,1),(1,-1)\bigr)$ est une base de $\R^{2}$ (vecteurs non
    colinéaires) : le théorème de détermination donne l'existence et
    l'unicité de $u$.
  \item $(x,y)=\dfrac{x+y}{2}(1,1)+\dfrac{x-y}{2}(1,-1)$, donc
    \[
    u(x,y)=\frac{x+y}{2}(2,0)+\frac{x-y}{2}(0,4)
    =\res{(x+y,\ 2x-2y)},
    \]
    et $u(1,0)=\res{(1,2)}$.
  \item $u(x,y)=0$ impose $x+y=0$ et $x-y=0$, donc $x=y=0$ :
    $\operatorname{Ker}u=\ens{0}$, $u$ injective, donc bijective
    (endomorphisme en dimension finie) : un automorphisme.
\end{questions}
\end{corrige}
\end{blocLivre}

\begin{blocLivre}{cor-12-14}
\begin{corrige}{14}{Le théorème du rang en pratique}
\demarche{Chercher les relations entre lignes ou colonnes, puis compter.}
\begin{questions}
  \item $L_3=L_1+L_2$ : le rang est au plus $2$ ; les deux premières
    colonnes $(1,0,1)$ et $(2,1,3)$ sont indépendantes :
    $\operatorname{rg}u=\res{2}$.
  \item Le système se réduit à $x+2y+t=0$, $y+z+t=0$ :
    $x=-2y-t$, $z=-y-t$, d'où la base de $\operatorname{Ker}u$ :
    \[
    \res{\bigl((-2,1,-1,0),\ (-1,0,-1,1)\bigr)},
    \]
    et $\operatorname{Im}u=\operatorname{Vect}\bigl((1,0,1),(2,1,3)\bigr)$.
    Vérification : $\dim\R^{4}=4=2+2$.
  \item Ni injective ($\operatorname{Ker}u\neq\ens{0}$), ni surjective
    ($\operatorname{rg}u=2<3$).
\end{questions}
\end{corrige}
\end{blocLivre}

\begin{blocLivre}{cor-12-15}
\begin{corrige}{15}{Projecteurs et symétries : les caractérisations}
\demarche{Tout découle de l'identité $x=(x-p(x))+p(x)$.}
\begin{questions}
  \item $(\mathrm{id}-p)^{2}=\mathrm{id}-2p+p^{2}=\mathrm{id}-p$. Si
    $x=p(x)$ alors $x\in\operatorname{Im}p$ ; réciproquement si $y=p(z)$,
    $p(y)=p^{2}(z)=p(z)=y$ : ainsi
    $\operatorname{Ker}(\mathrm{id}-p)=\ens{x\mid p(x)=x}=\operatorname{Im}p$.
  \item $x=\bigl(x-p(x)\bigr)+p(x)$ avec $p\bigl(x-p(x)\bigr)=0$ :
    $E=\operatorname{Ker}p+\operatorname{Im}p$. Et si
    $x\in\operatorname{Ker}p\cap\operatorname{Im}p$ : $x=p(x)=0$. Donc
    $E=\operatorname{Ker}p\oplus\operatorname{Im}p$.
  \item $s^{2}(x,y)=s(y,x)=(x,y)$ ;
    $\operatorname{Ker}(s-\mathrm{id})=\ens{(x,x)}$ est la diagonale,
    $\operatorname{Ker}(s+\mathrm{id})=\ens{(x,-x)}$ l'antidiagonale : $s$
    est la \res{\text{symétrie par rapport à la diagonale}} parallèlement à
    l'antidiagonale — et $p=\frac12(s+\mathrm{id})$ est le projecteur de
    l'exercice 7.
\end{questions}
\end{corrige}
\end{blocLivre}

\begin{blocLivre}{cor-12-16}
\begin{corrige}{16}{Hyperplans de $\R^{3}$}
\demarche{Un hyperplan, une forme linéaire : l'équation est unique à un
facteur près.}
\begin{questions}
  \item $\varphi$ est une forme linéaire non nulle : $H$ est un hyperplan,
    de dimension $2$. Les vecteurs $(2,1,0)$ et $(-1,0,1)$ vérifient
    $\varphi=0$ et ne sont pas colinéaires :
    $\res{\bigl((2,1,0),(-1,0,1)\bigr)}$ est une base de $H$.
  \item $\varphi(1,0,0)=1\neq0$ : la droite
    $D=\operatorname{Vect}\bigl((1,0,0)\bigr)$ n'est pas dans $H$, donc
    $\R^{3}=H\oplus D$. (Attention : $(1,1,1)$ serait un mauvais choix,
    car $\varphi(1,1,1)=0$.)
  \item $-2x+4y-2z=-2\,\varphi(x,y,z)$ : même noyau. En général, si
    $\operatorname{Ker}\varphi=\operatorname{Ker}\psi=H$, on choisit
    $v\notin H$ ; tout $x$ s'écrit $h+\lambda v$ et
    $\psi(x)=\lambda\psi(v)=\dfrac{\psi(v)}{\varphi(v)}\varphi(x)$ : deux
    équations d'un même hyperplan sont \res{\text{proportionnelles}}.
\end{questions}
\end{corrige}
\end{blocLivre}

\begin{blocLivre}{cor-12-17}
\begin{corrige}{17}{L'interpolation, vue par l'algèbre linéaire}
\demarche{Injectivité + égalité des dimensions : l'isomorphisme fait tout le
travail, sans résoudre le moindre système.}
\begin{questions}
  \item Linéarité claire. Si $\Phi(P)=0$, le polynôme $P$, de degré au plus
    $2$, possède trois racines $0,1,2$ : $P=0$ (chapitre 10). $\Phi$ est
    donc injective entre espaces de même dimension $3$ : c'est un
    isomorphisme.
  \item Pour tout $(y_0,y_1,y_2)$, il existe un unique $P\in\R_2[X]$ avec
    $\Phi(P)=(y_0,y_1,y_2)$ : existence et unicité du polynôme
    interpolateur, sans calcul.
  \item $\Phi(L_i)=e_i$ impose $L_i(a_j)=1$ si $i=j$, $0$ sinon :
    \[
    L_0=\frac{(X-1)(X-2)}{2},\qquad L_1=-X(X-2),\qquad
    L_2=\frac{X(X-1)}{2}.
    \]
    Enfin $\Phi(L_0+L_1+L_2)=(1,1,1)=\Phi(1)$ et $\Phi$ injective :
    $\res{L_0+L_1+L_2=1}$.
\end{questions}
\end{corrige}
\end{blocLivre}

\begin{blocLivre}{cor-12-18}
\begin{corrige}{18}{Solution particulière + solution homogène}
\demarche{Trois habits du même théorème : solutions = point particulier +
noyau.}
\begin{questions}
  \item Solution particulière $(1,0,2)$ ; système homogène : $x=y$,
    $z=-2x$, direction $\operatorname{Vect}\bigl((1,1,-2)\bigr)$ :
    \[
    \mathcal{S}=\res{(1,0,2)+\R\,(1,1,-2)},
    \]
    une droite affine de $\R^{3}$.
  \item Suite constante : $c=2c+3$ donne $c=-3$ ; solutions homogènes
    $\lambda\,2^{n}$ ; avec $u_0=1$ : $\lambda=4$, soit
    $\res{u_n=4\cdot2^{n}-3}$.
  \item Particulière $y=-1$ ; homogène $C\mathrm{e}^{x}$ :
    $\res{y=C\mathrm{e}^{x}-1}$, $C\in\R$.
\end{questions}
\end{corrige}
\end{blocLivre}

\begin{blocLivre}{cor-12-19}
\begin{corrige}{19}{Familles libres de fonctions}
\demarche{Faire parler l'infini : la croissance ou la non-dérivabilité
distingue les fonctions mieux que tout calcul.}
\begin{questions}
  \item Récurrence sur $n$ (cas $n=1$ clair). Soit
    $\lambda_1\mathrm{e}^{a_1x}+\dots+\lambda_n\mathrm{e}^{a_nx}=0$ pour
    tout $x$. En divisant par $\mathrm{e}^{a_nx}$ :
    \[
    \lambda_n=-\sum_{i<n}\lambda_i\,\mathrm{e}^{(a_i-a_n)x}
    \xrightarrow[x\to+\infty]{}0
    \]
    car $a_i-a_n<0$ : donc $\lambda_n=0$, et l'hypothèse de récurrence
    achève. La famille est \res{\text{libre}}.
  \item Soit $\sum_{j}\lambda_j\,|x-a_j|=0$ (somme finie, $a_j$ distincts).
    Si un $\lambda_j$ était non nul, la fonction nulle serait somme de
    fonctions toutes dérivables en $a_j$ \emph{sauf}
    $\lambda_j|x-a_j|$ : elle ne serait pas dérivable en $a_j$, absurde.
    Donc tous les $\lambda_j$ sont nuls : famille \res{\text{libre}}.
\end{questions}
\end{corrige}
\end{blocLivre}

\begin{blocLivre}{cor-12-20}
\begin{corrige}{20}{Endomorphismes nilpotents : premiers pas}
\demarche{Appliquer $u^{p-1}$, puis $u^{p-2}$… : chaque étape tue un
coefficient.}
\begin{questions}
  \item Soit $\lambda_0x+\lambda_1u(x)+\dots+\lambda_{p-1}u^{p-1}(x)=0$.
    En appliquant $u^{p-1}$ : $\lambda_0u^{p-1}(x)=0$ (tous les autres
    termes contiennent $u^{p}=0$), donc $\lambda_0=0$. En appliquant
    $u^{p-2}$ : $\lambda_1=0$, et ainsi de suite : la famille est libre.
    Elle compte $p$ vecteurs en dimension $n$ : $\res{p\leqslant n}$.
  \item Dériver baisse le degré de $1$ : $D^{n+1}=0$ sur $\R_n[X]$, et
    $D^{n}(X^{n})=n!\neq0$ : l'indice vaut exactement $\res{n+1}$ —
    cohérent avec a), puisque $\dim\R_n[X]=n+1$.
  \item $(\mathrm{id}-u)\bigl(\mathrm{id}+u+\dots+u^{p-1}\bigr)
    =\mathrm{id}-u^{p}=\mathrm{id}$, et les deux facteurs commutent :
    $\mathrm{id}-u$ est inversible, d'inverse la somme géométrique
    tronquée — exactement comme $\frac{1}{1-x}$ (chapitres 8 et 9).
\end{questions}
\end{corrige}
\end{blocLivre}

\begin{blocLivre}{cor-12-21}
\begin{corrige}{21}{Noyaux itérés}
\demarche{Une suite croissante de sous-espaces d'un espace de dimension
finie n'a pas le choix : elle stationne.}
\begin{questions}
  \item Si $u^{k}(x)=0$, alors $u^{k+1}(x)=u\bigl(u^{k}(x)\bigr)=0$ :
    l'inclusion est immédiate.
  \item Soit $x\in\operatorname{Ker}u^{k+2}$ : $u(x)\in\operatorname{Ker}u^{k+1}
    =\operatorname{Ker}u^{k}$, donc $u^{k+1}(x)=0$, c'est-à-dire
    $x\in\operatorname{Ker}u^{k+1}$ : l'égalité se propage de proche en
    proche à tous les rangs suivants.
  \item Tant que la suite ne stationne pas, les dimensions croissent
    strictement ; partant de $0$ et majorées par $n$, elles ne peuvent
    croître strictement plus de $n$ fois : la suite stationne au plus tard
    au rang $\res{n}$.
  \item Pour la dérivation sur $\R_3[X]$ :
    $\operatorname{Ker}D^{k}=\R_{k-1}[X]$ (les polynômes de degré
    $<k$), de dimensions $\res{0,1,2,3,4}$ pour $k=0,\dots,4$, puis
    stationnaire — la croissance est ici la plus lente possible.
\end{questions}
\end{corrige}
\end{blocLivre}

\begin{blocLivre}{cor-12-22}
\begin{corrige}{22}{Deux projecteurs qui commutent}
\demarche{La commutation permet de faire passer p au travers de q — chaque
inclusion s'écrit en une ligne.}
\begin{questions}
  \item $(pq)^{2}=p\,qp\,q=p\,pq\,q=p^{2}q^{2}=pq$ : $p\circ q$ est un
    projecteur.
  \item $\subseteq$ : $pq(x)=p\bigl(q(x)\bigr)\in\operatorname{Im}p$ et
    $pq(x)=qp(x)\in\operatorname{Im}q$. $\supseteq$ : si
    $y\in\operatorname{Im}p\cap\operatorname{Im}q$, alors $p(y)=y$ et
    $q(y)=y$ (exercice 15), donc $pq(y)=p(y)=y$ et
    $y\in\operatorname{Im}(pq)$. D'où
    $\res{\operatorname{Im}(pq)=\operatorname{Im}p\cap\operatorname{Im}q}$.
  \item $\supseteq$ : si $p(x)=0$, $pq(x)=qp(x)=0$ ; de même pour
    $\operatorname{Ker}q$, et une somme d'éléments des deux noyaux est
    encore dans $\operatorname{Ker}(pq)$. $\subseteq$ : si $pq(x)=0$,
    écrire $x=\bigl(x-q(x)\bigr)+q(x)$ : le premier terme vérifie
    $q\bigl(x-q(x)\bigr)=q(x)-q^{2}(x)=0$, le second
    $p\bigl(q(x)\bigr)=0$. D'où
    $\res{\operatorname{Ker}(pq)=\operatorname{Ker}p+\operatorname{Ker}q}$.
\end{questions}
\end{corrige}
\end{blocLivre}

\begin{blocLivre}{cor-12-23}
\begin{corrige}{23}{Les hyperplans de polynômes}
\demarche{« Racine » égale « facteur » : chaque condition d'annulation coûte
exactement une dimension.}
\begin{questions}
  \item $H_a$ est le noyau de la forme linéaire non nulle $P\mapsto P(a)$
    (elle vaut $1$ sur le polynôme constant $1$) : un hyperplan, de
    dimension $3$. Et $P(a)=0$ équivaut à $(X-a)\mid P$ (chapitre 10) :
    $H_a=(X-a)\,\R_2[X]$.
  \item $P\in H_0\cap H_1\cap H_2$ signifie que $X$, $X-1$, $X-2$ divisent
    $P$, de degré au plus $3$ : $P=\lambda\,X(X-1)(X-2)$. C'est la droite
    $\res{\operatorname{Vect}\bigl(X(X-1)(X-2)\bigr)}$.
  \item Le même argument donne $P=\prod_{i}(X-a_i)\,Q$ avec
    $\deg Q\leqslant3-(m+1)$ :
    \[
    \dim\bigl(H_{a_0}\cap\dots\cap H_{a_m}\bigr)=\res{3-m}.
    \]
    L'intersection de $m+1$ hyperplans atteint donc ici exactement la
    borne $n-(m+1)=3-m$ du cours : des conditions d'interpolation en des
    points distincts sont toujours « indépendantes ».
\end{questions}
\end{corrige}
\end{blocLivre}

\begin{blocLivre}{cor-12-24}
\begin{corrige}{24}{Quand $E=\operatorname{Ker}u\oplus\operatorname{Im}u$}
\demarche{Les inclusions sont toujours vraies ; le théorème du rang
transforme chaque égalité d'espaces en égalité de dimensions.}
\etape{b) équivaut à c) et d)}
On a toujours $\operatorname{Ker}u\subseteq\operatorname{Ker}u^{2}$ et
$\operatorname{Im}u^{2}\subseteq\operatorname{Im}u$. Le théorème du rang
appliqué à $u$ et à $u^{2}$ donne
$\dim\operatorname{Ker}u+\operatorname{rg}u=n
=\dim\operatorname{Ker}u^{2}+\operatorname{rg}u^{2}$ : l'égalité des noyaux
équivaut à celle des rangs, qui équivaut (par inclusion et égalité des
dimensions) à celle des images : (b) $\iff$ (d) $\iff$ (c).
\etape{a) implique b)}
Si $x\in\operatorname{Ker}u^{2}$, alors
$u(x)\in\operatorname{Ker}u\cap\operatorname{Im}u=\ens{0}$ : donc
$x\in\operatorname{Ker}u$.
\etape{b) implique a)}
Si $x\in\operatorname{Ker}u\cap\operatorname{Im}u$ : $x=u(y)$ et
$u(x)=0$, donc $y\in\operatorname{Ker}u^{2}=\operatorname{Ker}u$ et
$x=u(y)=0$ : l'intersection est nulle. Comme
$\dim\operatorname{Ker}u+\operatorname{rg}u=n$ (théorème du rang), la somme
est directe et vaut $E$.
\etape{Exemples}
Un projecteur vérifie $u^{2}=u$, donc (b) trivialement — et on retrouve
$E=\operatorname{Ker}p\oplus\operatorname{Im}p$ (exercice 15). Pour la
dérivation sur $\R_2[X]$ : $\operatorname{rg}D=2$ mais
$\operatorname{rg}D^{2}=1$, et
$\operatorname{Ker}D=\R_0[X]\subseteq\operatorname{Im}D=\R_1[X]$ : aucune
des quatre assertions n'est vérifiée.
\end{corrige}
\end{blocLivre}

% !TEX root = ../../main.tex
%% Corrigé du défi du chapitre 12 — vérifié par verifications/defis_11_15.py.

\begin{blocLivre}{cor-12-25a}
\begin{corrige}{25}{L'opérateur de différence et la formule de Newton}
\demarche{$\Delta$ fait baisser le degré de un : il est nilpotent, et la base
binomiale est exactement celle où il « décale » les vecteurs.}
\partieExo{A}{Un endomorphisme nilpotent}
\begin{questions}
  \item La linéarité est immédiate. Si $P=a_dX^{d}+a_{d-1}X^{d-1}+\cdots$,
    le binôme donne $P(X+1)=a_dX^{d}+(da_d+a_{d-1})X^{d-1}+\cdots$, donc
    $\Delta P=d\,a_dX^{d-1}+\cdots$ : $\Delta(E)\subseteq E$, et pour
    $d\geqslant1$, $\deg\Delta P=d-1$ avec \res{\text{coefficient dominant }d\,a_d}.
  \item $\Delta P=0$ si et seulement si $P$ est constant :
    $\res{\operatorname{Ker}\Delta=\R_0[X]}$. Donc
    $\operatorname{rg}\Delta=(n+1)-1=n$ ; comme
    $\operatorname{Im}\Delta\subseteq\R_{n-1}[X]$, de dimension $n$ :
    $\res{\operatorname{Im}\Delta=\R_{n-1}[X]}$. Chaque application de
    $\Delta$ fait baisser le degré, d'où $\Delta^{n+1}=0$ ; et
    $\Delta^{n}(X^{n})=n!\neq0$.
  \item $\Delta1=0$, $\Delta X=1$, $\Delta X^{2}=2X+1$,
    $\Delta X^{3}=3X^{2}+3X+1$ :
    \[
    \operatorname{Mat}(\Delta)=\begin{pmatrix}0&1&1&1\\0&0&2&3\\0&0&0&3\\0&0&0&0\end{pmatrix}
    \quad(\text{triangulaire stricte, de rang }3).
    \]
\end{questions}
\partieExo{B}{La base binomiale et la formule de Newton}
\begin{questions}[resume]
  \item $\deg H_k=k$ : famille échelonnée, donc base de $E$. Puis
    \[
    \Delta H_k=\frac{X(X-1)\cdots(X-k+2)}{k!}\bigl((X+1)-(X-k+1)\bigr)
    =\frac{X\cdots(X-k+2)}{(k-1)!}=\res{H_{k-1}} .
    \]
    Dans cette base, la matrice de $\Delta$ a des $1$ juste au-dessus de la
    diagonale et des $0$ ailleurs : $\Delta$ « décale » la base.
  \item Si $P=\sum_kc_kH_k$, alors $\Delta^{j}P=\sum_{k\geqslant j}c_kH_{k-j}$
    et, comme $H_0(0)=1$ et $H_m(0)=0$ pour $m\geqslant1$ :
    $\bigl(\Delta^{j}P\bigr)(0)=c_j$, d'où la formule. Pour $X^{3}$, la
    diagonale de la table donne $0,1,6,6$ :
    $\res{X^{3}=H_1+6H_2+6H_3}$.
  \item Pour $m<0$,
    $H_k(m)=(-1)^{k}\dfrac{(-m)(-m+1)\cdots(-m+k-1)}{k!}=(-1)^{k}\dbinom{k-m-1}{k}\in\Z$,
    et pour $m\geqslant0$, $H_k(m)=\binom mk$ : des coordonnées entières
    donnent des valeurs entières. Réciproquement,
    $\bigl(\Delta^{k}P\bigr)(0)=\sum_{j=0}^{k}(-1)^{k-j}\binom kj P(j)$ est
    entier si $P$ prend des valeurs entières. Exemple :
    $\frac{X(X+1)}{2}=\res{H_2+H_1}$.
\end{questions}
\end{corrige}
\end{blocLivre}

\begin{blocLivre}{cor-12-25b}
\begin{corrigeSuite}
\partieExo{C}{Sommation discrète}
\begin{questions}[start=7]
  \item $\sum_{k=0}^{m-1}P(k)=\sum_{k=0}^{m-1}\bigl(Q(k+1)-Q(k)\bigr)=Q(m)-Q(0)$
    (télescopage), exactement comme $\int_a^bf=F(b)-F(a)$.
  \item La table de $k^{2}$ ($0,1,4,9$ ; puis $1,3,5$ ; puis $2,2$) donne
    $X^{2}=H_1+2H_2$, et comme $\Delta H_{k+1}=H_k$, $Q=H_2+2H_3$
    convient. Ainsi
    \[
    \sum_{k=0}^{m-1}k^{2}=\binom m2+2\binom m3=\res{\frac{(m-1)m(2m-1)}{6}} .
    \]
    De même $X^{3}=H_1+6H_2+6H_3$ donne $Q=H_2+6H_3+6H_4$ et
    \[
    \sum_{k=0}^{m-1}k^{3}=\binom m2+6\binom m3+6\binom m4
    =\frac{m(m-1)}{4}\bigl(2+4(m-2)+(m-2)(m-3)\bigr)
    =\res{\Bigl(\frac{m(m-1)}{2}\Bigr)^{2}} ,
    \]
    soit $T_{m-1}^{2}$ : le théorème de Nicomaque, retrouvé par l'algèbre
    linéaire.
\end{questions}
\end{corrigeSuite}
\end{blocLivre}

